% chapters/11/section-02.tex
\section{Mengikuti Jalur}
\label{sec:mengikuti-jalur}

Rute hasil pencarian berupa rangkaian petak, sedangkan benda bergerak dalam
satuan piksel per detik. Bagian ini menjelaskan cara menjembatani keduanya.

\subsection{Bergerak Menuju Petak Berikutnya}

Pengikut jalur bekerja dengan aturan yang sederhana: lihat petak pertama pada
jalur, bergeraklah menuju titik tengahnya, dan apabila sudah cukup dekat,
buanglah petak itu dari rangkaian lalu ulangi dengan petak berikutnya.

\begin{kodekeluaran}
def ikuti_jalur(self, delta_seconds, kecepatan, dinding):
    while self.jalur:
        tujuan = pusat_petak(self.jalur[0])
        arah = tujuan - self.position
        if arah.length() <= REACH_DISTANCE:
            self.jalur.pop(0)
            continue
        self.position += arah.normalize() * kecepatan * delta_seconds
        break
\end{kodekeluaran}

Tiga hal perlu diperhatikan.

\begin{enumerate}
  \item \textbf{Perulangan \code{while}, bukan \code{if}.} Perulangan itu
        penting, sebab dalam satu bingkai NPC dapat melampaui lebih dari satu
        petak apabila kecepatannya tinggi atau lajunya rendah. Tanpa perulangan,
        petak-petak yang dilampaui itu tertinggal dan NPC berbalik untuk
        mencapainya.
  \item \textbf{Toleransi jarak \code{REACH_DISTANCE}.} Nilainya beberapa
        piksel, bukan nol. Tanpa toleransi itu, NPC tidak akan pernah tepat
        berada di titik tengah petak, sehingga petak itu tidak pernah dibuang
        dan NPC berhenti di tempatnya selamanya.
  \item \textbf{Jalur diperlakukan sebagai antrean.} Petak yang telah dilalui
        dibuang dari depan, sehingga petak berikutnya selalu berada di depan.
        Cara itu jauh lebih murah daripada menelusuri seluruh rangkaian setiap
        bingkai.
\end{enumerate}

\subsection{Jalur yang Menjadi Usang}

Ketika pemain berpindah, jalur yang sedang ditempuh menjadi usang: rangkaian
petaknya masih menuju tempat pemain berada sebelumnya. NPC akan terus
mengikutinya sampai selang hitung ulangnya habis, lalu memperbaiki arahnya.

Perilaku itu wajar, dan justru itulah pertukaran yang harus Anda nilai pada
bagian berikutnya. Memperpendek selang hitung ulang membuat NPC lebih tanggap
tetapi lebih berat dihitung, sedangkan memperpanjangnya menghemat perhitungan
tetapi membuat NPC sempat menempuh arah yang salah.

\begin{catatabox}
Pada permainan yang sesungguhnya, keusangan itu sering dikurangi dengan cara
lain yang lebih murah daripada menghitung ulang: jalur diperiksa apakah masih
memadai, misalnya dengan memeriksa apakah pemain masih berada di dekat ujung
jalur. Apabila ya, jalur itu cukup diperpanjang, bukan dihitung ulang
seluruhnya. Cara itu di luar cakupan bab ini, tetapi perlu Anda ketahui sebagai
arah pengembangan.
\end{catatabox}

\subsection{Gerak yang Tidak Selalu Mulus}

NPC bergerak menuju titik tengah setiap petak, sehingga
lintasannya berupa rangkaian garis lurus antar-petak, bukan garis halus. Pada
jalur yang berbelok tajam, geraknya terlihat menyudut.

Cara memperhalusnya adalah dengan menyingkirkan petak-petak yang tidak perlu,
yaitu petak-petak yang terletak pada satu garis lurus dengan petak sebelumnya
dan berikutnya. Cara itu disebut penghalusan jalur (\istilah{path smoothing}),
dan hasilnya adalah rangkaian titik belok saja, bukan seluruh petak.

\begin{tipsbox}
Sebelum menambahkan penghalusan jalur, perhatikan lebih dahulu apakah
memang diperlukan. Pada permainan dengan petak yang kecil dan jumlah NPC
yang sedikit, gerak menyudut itu hampir tidak terlihat. Menambahkan
penghalusan menambah perhitungan dan menambah satu hal yang dapat salah.
Kebiasaan mengukur lebih dahulu, baru memperbaiki, berlaku di sini.
\end{tipsbox}
